\section{De la autogestión por votación a la gobernanza por niveles}

Los primeros principios de Reddit se acercaban a «la plataforma no elimina; los usuarios deciden la visibilidad mediante votos». Cuando la escala era pequeña y la participación maliciosa limitada, esta autogestión reducía el poder editorial centralizado y además alentaba a las comunidades a formar su propia cultura. Al crecer la escala, la votación no podía atender el acoso entre comunidades, el contenido ilegal, la filtración de datos privados ni los grupos cuyo objetivo expreso es causar daño. La causa no es que los usuarios dejaran de ser confiables de repente, sino que el daño tiene externalidades: los afectados pueden no pertenecer a la comunidad que vota, y unos pocos coordinadores pueden manipular una mayoría local.

\subsection{Por qué la gobernanza tiene que organizarse por niveles}

Las reglas de la plataforma se encargan de la seguridad mínima, la legalidad y los límites entre comunidades; los moderadores de cada subreddit se encargan del tema y las normas locales; los votos de los usuarios se encargan del ordenamiento y la retroalimentación cotidianos. Cada uno de los tres niveles dispone de información distinta: la plataforma ve los patrones entre sitios y los riesgos legales, los moderadores entienden el contexto y la historia de la comunidad, y los usuarios comunes son quienes están más cerca del valor del contenido. Si cualquiera de los niveles concentra todo el poder, el resultado se distorsiona; el objetivo del sistema de gobernanza es que cada problema se atienda en el nivel que cuenta con información suficiente y cuyo alcance corresponde al del problema.

\begin{figure}[H]
\centering
\includegraphics[width=0.92\textwidth]{images/03-governance-stack.png}
\caption{Pila de gobernanza comunitaria: reparto de funciones entre la política de todo el sitio, la aplicación de normas por la plataforma, las reglas de cada subreddit y la participación de los usuarios.}
\figsource{Redibujado a partir de los subtítulos de la entrevista, 08:30--15:00, `images/03-governance-stack.png`.}
\end{figure}

\begin{knowledgebox}{Lectura de la figura: cuanto más alto el nivel, menos reglas pero límites más firmes}
La Site-wide policy se ocupa de los daños que atraviesan todo el sitio; debe ser relativamente estable, explicable y estar sujeta a la ley y a la responsabilidad pública. Las Subreddit rules pueden ser más detalladas y permiten que cada comunidad elija el formato de las publicaciones, los temas y el estilo de interacción. Voting es una señal débil y de alta frecuencia: sirve para ordenar, no para dictaminar sobre todos los problemas de seguridad. Las fallas de gobernanza suelen ocurrir por un desajuste de niveles: dejar que la votación resuelva la exposición de datos personales, o que la sede de la plataforma decida el tono de cada comunidad pequeña.
\end{knowledgebox}

\begin{importantbox}{Términos de primera aparición: moderation y doxing}
Moderation es la gobernanza del contenido y de la comunidad: incluye establecer reglas, revisar denuncias, restringir cuentas, atender apelaciones y apoyar a los moderadores. Doxing es publicar sin consentimiento información sensible de una persona, como su identidad, su domicilio o sus datos de contacto, lo que puede convertir un conflicto en línea en un daño en el mundo real. No puede decidirse solo por popularidad: requiere restricciones a nivel de todo el sitio y una respuesta rápida.
\end{importantbox}

\subsection{``Specifically vague'': por qué las reglas no pueden redactarse como una lista exhaustiva}

Los participantes maliciosos buscan resquicios en la letra de las reglas. Si una regla solo enumera palabras o conductas fijas, el atacante puede cambiar la ortografía, dividir sus acciones u organizar a otros para causar el mismo daño. Una regla demasiado amplia, en cambio, impide que los usuarios comunes anticipen dónde está el límite, y el equipo encargado de aplicarla tiende a actuar de forma inconsistente. Huffman usa ``specifically vague'' para describir el intervalo de trabajo: explicar con claridad el daño que se quiere evitar y el principio que lo sustenta, y conservar a la vez margen para juzgar, según el contexto, las conductas de evasión equivalentes.

\begin{figure}[H]
\centering
\includegraphics[width=0.92\textwidth]{images/04-policy-specificity.png}
\caption{El diseño de políticas necesita encontrar un intervalo de trabajo entre la previsibilidad y la resistencia a la evasión literal.}
\figsource{Redibujado a partir de los subtítulos de la entrevista, 13:40--16:00, `images/04-policy-specificity.png`.}
\end{figure}

\begin{knowledgebox}{Lectura de la figura: el margen de ambigüedad debe estar acotado por procedimientos}
La zona intermedia no autoriza a quien aplica las normas a decidir arbitrariamente. Necesita ejemplos, precedentes, revisión entre equipos, apelaciones, informes de transparencia y revisiones periódicas que la acoten. Cuanto más se dependa del juicio contextual, más hay que registrar por qué casos similares recibieron un trato distinto. El texto de la política expresa el objetivo; la justicia procedimental evita que el objetivo quede capturado por sesgos personales.
\end{knowledgebox}

\begin{warningbox}{Nota de clase: cada palabra proviene de un costo}
El ponente subraya que las frases de una política de contenido madura suelen ser el resultado condensado de años de incidentes, controversias y casos límite. Al leer diez reglas, no basta con evaluar si la redacción es elegante: también hay que preguntar si existen las herramientas correspondientes de detección, revisión humana, notificación, apelación y restablecimiento. Una regla sin un sistema que la aplique solo traslada la responsabilidad a los equipos de primera línea.
\end{warningbox}

\subsection{La aplicación de normas pasa de un interruptor binario a una escalera}

Las primeras herramientas se reducían a la suspensión permanente o a no hacer nada. Una elección binaria lleva al equipo a aplicar sanciones excesivas o, por temor a una pérdida irreversible, a no actuar en absoluto. Los sistemas más maduros forman una escalera con la explicación, la advertencia, la restricción de corto plazo, el subreddit timeout, la suspensión temporal y el retiro permanente. La intensidad de la intervención aumenta según el daño, el número de reincidencias, la disposición a cooperar y la posibilidad de recuperación, y se ofrece una vía de reversión para los errores de juicio.

\begin{figure}[H]
\centering
\includegraphics[width=0.94\textwidth]{images/05-enforcement-ladder.png}
\caption{Aplicación gradual de normas: de la explicación y la advertencia se escala progresivamente hasta la restricción temporal y el retiro permanente.}
\figsource{Redibujado a partir de los subtítulos de la entrevista, 16:00--20:00, `images/05-enforcement-ladder.png`.}
\end{figure}

\begin{importantbox}{Lectura de la figura: la precisión de las herramientas determina la precisión de la gobernanza}
Si el sistema solo cuenta con el ban permanente, el equipo de políticas no puede expresar diferencias aunque las entienda. El timeout permite a la plataforma pausar una comunidad que se salió de control, dialogar con los moderadores y conservar la posibilidad de restablecerla; la suspensión escalonada da a quien infringe de forma ocasional la oportunidad de corregirse y, al mismo tiempo, acumula evidencia sobre la conducta maliciosa persistente. Por ello, la calidad de la gobernanza resulta de la combinación de políticas, datos y herramientas de producto.
\end{importantbox}

\begin{table}[H]
\centering
\begin{tabular}{L{0.18\textwidth}L{0.27\textwidth}L{0.43\textwidth}}
\toprule
Intervención & Situación aplicable & Procedimiento que debe acompañarla \\
\midrule
Explicación/recordatorio & Primera vez, bajo riesgo, regla poco clara & Señalar la regla y el contenido concretos \\
Advertencia & Infracción confirmada pero corregible & Conservar el registro y explicar las consecuencias de reincidir \\
Restricción temporal & Infracción reiterada o riesgo de corto plazo & Precisar plazo, alcance y condiciones de restablecimiento \\
Timeout de la comunidad & Los moderadores perdieron el control o el conflicto escala & Dialogar con los moderadores, proteger el contenido histórico, revisar antes de reabrir \\
Retiro permanente & Daño grave, evasión maliciosa o fallas persistentes & Revisión de alto nivel, conservación de evidencia, apelación e informe de transparencia \\
\bottomrule
\end{tabular}
\caption{La escalera de aplicación de normas reúne en una misma tabla el daño, la reversibilidad y las garantías procedimentales.}
\end{table}

\subsection{Resumen de la sección}

La votación sirve para ordenar, pero no basta para asumir toda la responsabilidad sobre la seguridad y la legalidad. La evolución de la gobernanza de Reddit organizó por niveles la política de todo el sitio, la aplicación de normas por la plataforma, las reglas de cada comunidad y la participación de los usuarios, y recurrió a ``specifically vague'' y a herramientas graduales para atender la evasión maliciosa y los casos límite. La siguiente sección trata otro límite: si el hecho de que un contenido sea de lectura pública significa que cualquier entidad puede extraerlo sin límite y explotarlo comercialmente.
